% Origen: sección "Apéndice" de Secciones/PI - Analizador Multicanal.md
\chapter{Complementos del analizador multicanal}
\label{C:ap-mca}

% Figuras de este apéndice + las compartidas.
\graphicspath{{figs/apendices/}{figs/comunes/}}

\section{Posibles detectores empleados}
\label{S:ap-detectores}

En la Tabla~\ref{t:ap-detectores} se presentan algunas de las familias de
detectores con las que se emplea un MCA, la partícula que detectan y los órdenes
de magnitud de sus tiempos característicos, su rango de energía y su tasa de
conteo típica~\cite{Knoll2010}.

\begin{table}[H]
\centering
\caption[Familias de detectores y sus parámetros característicos]{%
Familias de detectores con sus tiempos característicos, rangos de energía y
tasas de conteo típicas.\label{t:ap-detectores}}
\footnotesize
\renewcommand{\arraystretch}{1.3}
\setlength{\tabcolsep}{3pt}
\newcolumntype{Y}{>{\raggedright\arraybackslash}X}
\begin{tabularx}{\textwidth}{|>{\raggedright\arraybackslash}p{0.15\textwidth}|Y|Y|Y|Y|Y|}
\hline
\encabezado \textbf{Detector} & \textbf{Partícula detectada} & \textbf{Tiempo de subida}
  & \textbf{Tiempo de decaimiento / muerto} & \textbf{Rango de energía típico}
  & \textbf{Tasa de conteo típica (orden de magnitud)} \\ \hline
Cámara de ionización
  & $\alpha$, $\beta$, $\gamma$, rayos X, n (con convertidor)
  & \unit{\micro\second} (electrones) -- \unit{\milli\second} (iones)
  & \unit{\milli\second} (cola iónica)
  & \unit{\kilo\electronvolt} -- decenas de \unit{\mega\electronvolt}
  & $10^2$--$10^3$ eventos/s (modo pulso; en modo corriente no aplica) \\ \hline
Contador proporcional
  & $\alpha$, $\beta$, rayos X blandos, n térmicos (BF$_3$, $^3$He)
  & $\sim$\unit{\nano\second} -- cientos de \unit{\nano\second} (electrónica)
  & \unit{\milli\second} si no hay \emph{pulse shaping} CR-RC
  & \unit{\kilo\electronvolt} -- pocos \unit{\mega\electronvolt}
  & $10^4$--$10^5$ eventos/s \\ \hline
Geiger-Müller
  & Cualquier radiación ionizante (sin discriminación)
  & $\sim$\unit{\micro\second}
  & \qtyrange{50}{200}{\micro\second} (tiempo muerto)
  & No espectroscópico
  & $10^3$--$10^4$ eventos/s (saturación por tiempo muerto arriba de esto) \\ \hline
\textbf{Centellador orgánico (plástico, antraceno)}
  & $\beta$, n rápidos (retroceso de protones), $\gamma$ (timing)
  & $\sim$\qtyrange{1}{3}{\nano\second} (plástico), $\sim$\qty{30}{\nano\second} (antraceno)
  & \unit{\nano\second} (decaimiento rápido)
  & \unit{\kilo\electronvolt} -- \unit{\mega\electronvolt}
  & $10^6$--$10^7$ eventos/s \\ \hline
NaI(Tl)
  & $\gamma$, rayos X
  & $\sim$decenas de \unit{\nano\second}
  & $\sim$\qty{230}{\nano\second}
  & \unit{\kilo\electronvolt} -- varios \unit{\mega\electronvolt}
  & $10^5$--$10^6$ eventos/s \\ \hline
CsI(Tl)
  & $\gamma$, rayos X
  & $\sim$decenas de \unit{\nano\second}
  & $\sim$\qty{1000}{\nano\second}
  & \unit{\kilo\electronvolt} -- varios \unit{\mega\electronvolt}
  & $10^4$--$10^5$ eventos/s \\ \hline
LaBr$_3$(Ce)
  & $\gamma$, rayos X
  & $\sim$\unit{\nano\second}
  & \qtyrange{16}{25}{\nano\second}
  & \unit{\kilo\electronvolt} -- varios \unit{\mega\electronvolt}
  & $10^6$--$10^7$ eventos/s \\ \hline
Si (Si(Li), SDD, barrera superficial)
  & $\alpha$, $\beta$, rayos X blandos
  & $\sim$\qtyrange{10}{50}{\nano\second}
  & \unit{\nano\second} -- \unit{\micro\second} (según \emph{shaping})
  & \unit{\kilo\electronvolt} -- pocos \unit{\mega\electronvolt}
  & $10^5$--$10^6$ eventos/s \\ \hline
HPGe
  & $\gamma$ (principalmente)
  & decenas -- cientos de \unit{\nano\second}
  & \unit{\micro\second} (limitado por \emph{shaping} para resolución óptima)
  & \unit{\kilo\electronvolt} -- varios \unit{\mega\electronvolt}
  & $10^4$--$10^5$ eventos/s \\ \hline
CdTe / CZT
  & $\gamma$, rayos X
  & $\sim$decenas -- cientos de \unit{\nano\second}
  & \unit{\nano\second} -- \unit{\micro\second}
  & \unit{\kilo\electronvolt} -- $\sim$\qty{1}{\mega\electronvolt}
  & $10^4$--$10^5$ eventos/s \\ \hline
\end{tabularx}
\end{table}

\section{Aplicaciones del MCA 2D}
\label{S:ap-mca2d}

En la Tabla~\ref{t:ap-mca2d} se presentan las aplicaciones más importantes de
los MCA 2D intracanal.

\begin{table}[H]
\centering
\caption[Principales aplicaciones del MCA 2D intracanal]{%
Principales aplicaciones del MCA 2D intracanal.\label{t:ap-mca2d}}
\small
\renewcommand{\arraystretch}{1.3}
\begin{tabularx}{\textwidth}{|>{\raggedright\arraybackslash}p{0.22\textwidth}|>{\raggedright\arraybackslash}X|>{\raggedright\arraybackslash}p{0.27\textwidth}|}
\hline
\encabezado \textbf{Aplicación} & \textbf{Qué se explota} & \textbf{Requerimientos} \\ \hline
\textbf{Discriminación n/$\gamma$ (PSD)}
  & el centelleador orgánico (EJ-301, BC-501A, estilbeno) da \textbf{colas de
    decaimiento distintas} para protón de retroceso y electrón; los neutrones
    forman un lóbulo propio en el mapa
  & Eje de \emph{fall time} \\ \hline
\textbf{Identificación $\alpha/\beta/\gamma$}
  & Para phoswich, ZnS(Ag)+plástico, CsI(Tl): \textbf{componente rápida/lenta
    distinta por partícula}
  & soportado conceptualmente por el mismo eje \\ \hline
\textbf{Rechazo de colección incompleta de carga}
  & en CdZnTe/HPGe los eventos de mala forma (tiempo de subida largo, eventos
    de superficie) son los que \textbf{fabrican la cola de baja energía} del
    fotopico
  & Eje de \emph{rise time}. \\ \hline
\textbf{Diagnóstico de la propia cadena}
  & deriva de línea base, saturación y déficit balístico mueven o estiran el
    lóbulo; en 1D son invisibles
  & un defecto de la cadena que en 1D pasa inadvertido aparece en el mapa 2D
    como un «pico fantasma» (ver
    Sec.~\ref{S:intro-hw-formal}) \\ \hline
\end{tabularx}
\end{table}

% Las tablas son más largas que el texto: sin el \clearpage flotan detrás de
% las referencias.
\clearpage
% Los apéndices también son \chapter, así que tienen su propia bibliografía.
\bibliografiacapitulo

%%% Local Variables:
%%% mode: latex
%%% TeX-master: "template"
%%% End:
